\section{Gnome Sort: a Data-Dependent Number of Iterations}\label{sec:gnome}

\subsection{The program}
\begin{lstlisting}
method main(a: array<int, int>, n: int) returns (b: array<int, int>)
  requires 0 <= n
  ensures sorted(b, 0, n)
  ensures perm(b, a)
{
  var i: int;
  var t: int;
  b := a;
  i := 1;
  while i < n
  {
    if i == 0 || b[i - 1] <= b[i] {
      i := i + 1;
    } else {
      t := b[i];
      b[i] := b[i - 1];
      b[i - 1] := t;
      i := i - 1;
    }
  }
}
\end{lstlisting}
The position moves up when the adjacent elements are in order and down with a swap otherwise. The number of
iterations is not a counter and is unrelated to $n$: the exit $n \le i$ bounds it from below only. The file
was the folder's negative example, "open, as expected", until it was verified with two lemmas.

\subsection{The trace equation system}
\begin{align}
  K &= \cnt(0, 1, n') \qquad i = \cur(I) \qquad b = \cur(B) \qquad n \le i \label{eq:gn0}\\
  I &= \tr{1}{\bl{K}, \bl{B}}{i, k, b \mapsto \upd{G}{i + 1}\ \ \upd{\neg G}{i - 1}} \label{eq:gn1}\\
  B &= \tr{a}{\bl{K}, \bl{I}}{b, k, i \mapsto \upd{G}{b}\ \ \upd{\neg G}{b[i := b[i - 1]][i - 1 := b[i]]}} \label{eq:gn2}\\
  &\qquad\text{where } G \text{ abbreviates } i = 0 \vee b[i - 1] \le b[i] \nonumber\\
  &\al{\bl{K}, \bl{I}}{k, i \mapsto i < n} \label{eq:gn3}
\end{align}
As in the merge and the lower bound, two traces over each other with a data-dependent guard.

\subsection{The lemmas}
The invariant of gnome sort is $\sorted(b, 0, i)$: a step up extends the sorted prefix by the element that
is in order with its predecessor, a swap down keeps the prefix below the position sorted, and at the exit
$n \le i$ the prefix covers the range. The lemma states this on the whole atom of the postcondition
(Section~\ref{sec:merge}), with the guard and the swap positions higher-order and pinned, since the literal
$i - 1$ is not matched:
\begin{align*}
  &\lemma{sorted\_gnome}:\quad Bv = \tr{a}{\bl{Kv}, \bl{Iv_2}}{b, k, i \mapsto \upd{G(i, b)}{b}\ \ \upd{\neg G(i, b)}{b[Q(i) := b[P(i)]][P(i) := b[Q(i)]]}}\\
  &\quad\wedge\ Iv_2 = Iv\\
  &\quad\wedge\ Iv_2 = \tr{1}{\bl{Kv_2}, \bl{Bv_2}}{i, k, b \mapsto \upd{G(i, b)}{i + u}\ \ \upd{\neg G(i, b)}{i + d}}\\
  &\quad\wedge\ Kv_2 = Kv \wedge Bv_2 = Bv \wedge u = 1 \wedge d = -1\\
  &\quad\wedge\ (\forall i, b.\ G(i, b) \Leftrightarrow (i = 0 \vee b[i - 1] \le b[i]))\ \wedge\ (\forall i.\ P(i) = i - 1)\ \wedge\ (\forall i.\ Q(i) = i)\\
  &\quad\wedge\ x = \cur(Iv)\ \wedge\ n \le x\\
  &\quad\Longrightarrow\ \sorted(\cur(\tr{a}{\bl{Kv}, \bl{Iv}}{b, k, i \mapsto \upd{G(i, b)}{b}\ \ \upd{\neg G(i, b)}{b[Q(i) := b[P(i)]][P(i) := b[Q(i)]]}}), 0, n)\
    \Longleftrightarrow\ \mathsf{true}
\end{align*}
The left-hand side matches $\sorted(b, 0, n)$ through $b = \cur(B)$ and \eqref{eq:gn2}, binding $Kv := K$,
$Iv := I$, $G$, $P := \lambda i.\ i - 1$, $Q := \lambda i.\ i$ by abstraction. The conditions follow the
discipline of Section~\ref{sec:fd} and the lower bound: the trace of the array is tied first, open in $Bv$
and $Iv_2$, matched against \eqref{eq:gn2}, and $Iv_2 = Iv$ is decided; the trace of the position is tied
second, with fresh $Kv_2$, $Bv_2$ and the steps $u$, $d$ as variables (a literal $-1$ in a pattern is not
matched), matched against \eqref{eq:gn1}, and identified; the pins are decided by beta reduction; and the
exit is read from the facts through $x = \cur(Iv)$, matched against $i = \cur(I)$, then $n \le x$ decided.
No condition mentions the left-hand side, so nothing re-enters. For the permutation, \lemma{perm\_gnome}
states that every step of such a trace is the identity or a swap, for any guard and positions, hence
$\perm(\cur(\ldots), c) \Leftrightarrow \perm(a, c)$, and \lemma{perm\_refl} closes it. A wrong postcondition
$\sorted(b, 0, n + 1)$ leaves the goal.

\subsection{The idea}
A data-dependent number of iterations is no obstacle when the property can be stated on the final values
alone: the lemma does not describe how many steps were taken, only that the invariant holds after any number
of them and that the exit turns it into the postcondition. This is the same shape as the merge, where the
pointers are traces over each other, and it shows that the "negative example" of the paper's loop shapes is
negative for the characterization of the traces, not for the verification of the postcondition.
